% GENERATED FILE. Edit canonical lessons, then run npm run generate:books.
% canonical-source-hash: fnv1a64:bcf6965de3964ce8
% canonical-lessons: ZH-W07-zao, ZH-W07-shang, ZH-C07-zaoshang, ZH-C07-zaoshanghao, ZH-C07-practice

\chapter{Good Morning, and a Rule That Declines to Fire}
\label{ch:zh-zaoshang}
\hlchaptermodality{7}

\begin{chapteropening}
\textbf{By the end of this chapter:} \emph{I can write \zh{早} and \zh{上}, greet someone with \zh{早上好}, hear a neutral-tone syllable as the absence of a tone mark rather than a mistake, and say why the third-tone rule changes \zh{你好} but leaves \zh{早上好} alone.}

The first exchange where the reader chooses between two correct greetings rather than reciting the only one available -- and the first place all three tone behaviours in the track line up side by side, one moving and two declining to.
\end{chapteropening}

\section[zǎo]{\zh{早} — sun on top, one more time}

\label{lesson:ZH-W07-zao}



\begin{quote}
Write \zh{日}. Four strokes. Middle bar before the bottom bar.
\end{quote}

\subsection*{Script — the shape}
\begin{quote}
\zh{早}
\end{quote}

\textbf{Six strokes}, and you can already guess how it starts.

\textbf{\zh{早} is \zh{日} on top of \zh{十}.} You have written the sun before. The piece underneath is \textbf{\zh{十}}, which on its own is the character for \textbf{ten} — you have not been taught it and you do not need it today. Treat it as two strokes: one across, one down through it.

1-4. \textbf{\zh{日}, complete.} Left side; top and right in one turning stroke; the middle bar; the bottom bar closing it. 5. \textbf{\zh{十}'s horizontal}, left to right 6. \textbf{\zh{十}'s vertical}, coming down \emph{through} that horizontal

Five pen lifts.

\textbf{Predict the rule before you read it: does \zh{日} get finished first?} It does — the same habit that made \zh{是} manageable, and by now it should be doing the work for you rather than the instruction.

A word about what \zh{早} means, offered as a memory hook and labelled as one: the sun sitting above a short line is easy to read as \textbf{the sun just up over the horizon}, and \zh{早} means \textbf{early}. That is a picture worth keeping. It is not a proof, and the honest position is that the lower piece's history is disputed.

\subsection*{Your turn}
\begin{itemize}
\raggedright
  \item \emph{Write it:} \zh{日}, then \zh{早} — and stop after four strokes to check the sun is closed
  \item \emph{Say it:} \textbf{zǎo} — low and dipping
\end{itemize}

\begin{tcolorbox}[breakable,colback=teal!4,colframe=teal!35!black,title={Before you move on}]
Which way does the last stroke go? (\textbf{Down, through the horizontal}.)

Source: \href{https://www.unicode.org/charts/PDF/U4E00.pdf}{Unicode CJK Unified Ideographs chart}; stroke order per the PRC standard \emph{Xiandai Hanyu Tongyongzi Bishun Guifan} (1997).
\end{tcolorbox}
\section[shàng]{\zh{上} — and a rule that does not apply}

\label{lesson:ZH-W07-shang}



\begin{quote}
Write \zh{早}. Sun first, then the two strokes below.
\end{quote}

\subsection*{Script — the shape}
\begin{quote}
\zh{上}
\end{quote}

\textbf{Three strokes.} A vertical, a short horizontal to its right, and a long horizontal underneath.

Now guess the order, and expect to be wrong:

1. the \textbf{central vertical}, from the top down 2. the \textbf{short horizontal}, left to right, starting at the vertical 3. the \textbf{long base horizontal}, left to right

Two pen lifts.

\textbf{The vertical comes first.} Every character so far has trained you toward \emph{top before bottom, left before right}, and here the top stroke is not the one you start with. This is worth meeting on a three-stroke character, because it tells you something true about the system: the habits are \textbf{strong tendencies with named exceptions}, not laws. \zh{上} is one of the exceptions, and it is common enough that you will see it constantly.

The other thing to notice: \textbf{short horizontal before long horizontal.} The base bar is the widest stroke and it goes last, exactly as the closing bar did in \zh{再}. When two horizontals differ in length, the long one is usually the finisher.

\subsection*{Your turn}
\begin{itemize}
\raggedright
  \item \emph{Write it:} \zh{上} — vertical first, and resist starting at the top left
  \item \emph{Write it:} \zh{上} again, and say \textquotedblleft{}short, then long\textquotedblright{} as you draw the two bars
  \item \emph{Say it:} \textbf{shàng} — falling
\end{itemize}

\begin{tcolorbox}[breakable,colback=teal!4,colframe=teal!35!black,title={Before you move on}]
Which of the two horizontals is drawn last? (\textbf{The long one}.)

Source: \href{https://www.unicode.org/charts/PDF/U4E00.pdf}{Unicode CJK Unified Ideographs chart}; stroke order per the PRC standard \emph{Xiandai Hanyu Tongyongzi Bishun Guifan} (1997).
\end{tcolorbox}
\section[zǎoshang]{\zh{早上} — morning, and a syllable that goes light}

\label{lesson:ZH-C07-zaoshang}



\begin{quote}
Write \zh{早}. Six strokes, sun finished first.
\end{quote}

\begin{cousinweb}
\begin{quote}
\zh{早上}
\end{quote}

\textbf{\zh{早上}} \emph{zǎoshang} — \textbf{morning}.

The parts add up again, though less tidily than \zh{再见} did. \textbf{\zh{早}} is \emph{early}; \textbf{\zh{上}} on its own means \emph{up} or \emph{on top}. Early-up, the early part of the day. Take that as a hook rather than a definition: what you need is that the two characters you just wrote make the word for morning.
\end{cousinweb}

\begin{sounds}
The pinyin for this word is \emph{zǎoshang} — and \textbf{\zh{上} has lost its tone mark.}

On its own \zh{上} is \emph{shàng}, fourth tone, falling — you practised it that way. Inside this word it goes \textbf{neutral}: short, light, and pitched by whatever came before it. Neutral syllables are written with \textbf{no mark at all}, and that absence is the information.

This is a fifth thing a syllable can do, alongside the four tones. It is not a tone that moves — you have met two of those — it is a syllable \textbf{giving up its tone} because it is unstressed. {[}REPEAT x2{]}

Say it: \emph{zǎo} low and dipping, then \emph{shang} light and quick, riding on the end. The commonest mistake is to give the second syllable a full falling tone out of habit, which makes the word sound like two words.
\end{sounds}

\subsection*{Your turn}
Say these aloud:

\begin{itemize}
\raggedright
  \item \textbf{shàng} alone, falling — then \textbf{zǎoshang}, with the second syllable light
  \item \textbf{zǎoshang} twice more, and say aloud which syllable carries no tone
\end{itemize}

\begin{tcolorbox}[breakable,colback=teal!4,colframe=teal!35!black,title={Before you move on}]
Does \zh{上} keep its falling tone when it stands alone? (\textbf{Yes} — it only goes light inside this word.)

Source: \href{https://www.unicode.org/charts/PDF/U4E00.pdf}{Unicode CJK Unified Ideographs chart}
\end{tcolorbox}
\section[zǎoshang hǎo]{\zh{早上好} — the greeting was a pattern all along}

\label{lesson:ZH-C07-zaoshanghao}



\begin{quote}
Say \textbf{zǎoshang}. Second syllable light.
\end{quote}

\begin{cousinweb}
\begin{quote}
\zh{早上好}
\end{quote}

\textbf{\zh{早上好}} \emph{zǎoshang hǎo} — \textbf{good morning}.

It ends in \textbf{\zh{好}} — the same character, the same word, that you met in the first greeting in this book.

That is the payoff, and it is worth being explicit about. \textbf{\zh{你好}} was never a phrase to swallow whole. It is \zh{你} \emph{you} followed by \zh{好} \emph{good}, and the second half is reusable:

\begin{quote}
\zh{你} + \zh{好} $\to$ \textbf{\zh{你好}}
\end{quote}

>

\begin{quote}
\zh{早上} + \zh{好} $\to$ \textbf{\zh{早上好}}
\end{quote}

\textbf{Put a person or a time in front of \zh{好} and you have a greeting.} You have not learned two greetings; you have learned a \textbf{frame} and filled it twice. Every one of the other times of day works the same way, and you will be able to build those greetings the moment you have the words for those times.
\end{cousinweb}

\begin{sounds}
You know a rule about two third tones meeting: the first one rises, which is why \zh{你好} is said \emph{ní hǎo}.

\textbf{\zh{早}} is third tone. \textbf{\zh{好}} is third tone. So does \zh{早上好} become \emph{záoshang hǎo}?

\textbf{No — and the reason is the whole point of a rule having conditions.} The rule needs the two third tones to be \textbf{next to each other}. Here they are not: \zh{上} sits between them. A light, toneless syllable is still a syllable, and it keeps the two dipping tones apart.

So \emph{zǎoshang hǎo}: \zh{早} low, \zh{上} light, \zh{好} low and dipping. Nothing rises.
\end{sounds}

\subsection*{Your turn}
Say these aloud:

\begin{itemize}
\raggedright
  \item \textbf{nǐ hǎo}, then \textbf{zǎoshang hǎo} — and say which one moved
  \item \textbf{zǎoshang}, then \textbf{zǎoshang hǎo} — the greeting is the word plus \zh{好}
\end{itemize}

\begin{tcolorbox}[breakable,colback=teal!4,colframe=teal!35!black,title={Before you move on}]
Why does the third-tone rule not fire in this greeting? (\textbf{The two third tones are not next to each other}.)

Source: \href{https://www.unicode.org/charts/PDF/U4E00.pdf}{Unicode CJK Unified Ideographs chart}
\end{tcolorbox}
\section[Practice]{Practice — greeting the time of day}

\label{lesson:ZH-C07-practice}



\begin{quote}
Say \textbf{zǎoshang hǎo}. Low, light, low.
\end{quote}

\subsection*{The exchange}
\begin{quote}
\textbf{A:} \zh{早上好}
\end{quote}

>

\begin{quote}
\emph{zǎoshang hǎo} — good morning.
\end{quote}

>

\begin{quote}
\textbf{B:} \zh{早上好}
\end{quote}

>

\begin{quote}
\emph{zǎoshang hǎo} — good morning.
\end{quote}

>

\begin{quote}
\textbf{A:} \zh{再见}
\end{quote}

>

\begin{quote}
\emph{zàijiàn} — goodbye.
\end{quote}

>

\begin{quote}
\textbf{B:} \zh{再见}
\end{quote}

>

\begin{quote}
\emph{zàijiàn} — goodbye.
\end{quote}

Short, and for the first time you had a \textbf{choice} to make. \zh{你好} would also have been correct. \zh{早上好} is the one that fits the hour, and choosing between them is a thing a speaker does, not a thing a phrasebook does for you.

\begin{grammarlens}[title={What you've built — three tones, three behaviours}]
Line the greetings up and every tone behaviour you have been taught is visible at once:

\begin{itemize}
\raggedright
  \item \textbf{\zh{你好}} — two third tones, side by side. The first \textbf{rises}: \emph{ní hǎo}.
  \item \textbf{\zh{早上好}} — two third tones, held apart by a light syllable. \textbf{Neither moves}.
  \item \textbf{\zh{再见}} — two falling tones. \textbf{Neither moves.}
\end{itemize}

That middle line is the one that repays study. It is the same rule as the first line, and it does not fire, because the condition it needs is not met. \textbf{A rule you can only apply is half a rule; a rule you can also decline to apply is a whole one.}

And none of it shows on the page. Every character in all three is written exactly as it would be alone.
\end{grammarlens}

\subsection*{Your turn}
\begin{itemize}
\raggedright
  \item \emph{Say it:} \textbf{nǐ hǎo}, then \textbf{zǎoshang hǎo} — and say which one moved
  \item \emph{Say it:} the whole exchange, both parts
  \item \emph{Write it:} \zh{早上好}, then \zh{再见}
\end{itemize}

\begin{tcolorbox}[breakable,colback=teal!4,colframe=teal!35!black,title={Before you move on}]
Which of the three expressions above has a tone that moves? (\textbf{Only \zh{你好}}.)

Source: \href{https://www.unicode.org/charts/PDF/U4E00.pdf}{Unicode CJK Unified Ideographs chart}
\end{tcolorbox}
